\documentclass{article}
\usepackage{amsmath,amssymb,amsthm}
\usepackage{algorithm}
\usepackage{algorithmic}
\usepackage{bm}
\usepackage{enumitem}
\usepackage{hyperref}
\newtheorem{theorem}{Theorem}
\newtheorem{lemma}{Lemma}
\newtheorem{assumption}{Assumption}
\newcommand{\prox}{\operatorname{prox}}
\newcommand{\grad}{\nabla}
\begin{document}

\section*{Algorithm Name}
\textbf{Proximal Gradient Method with Gradient Momentum (PGM‑M)}  

The method solves the composite convex problem  

\[
\min_{x\in\mathbb{R}^d}\;F(x):=f(x)+g(x),
\]

where $f$ is smooth and convex and $g$ is convex (possibly nonsmooth) with an easy proximal operator.

\section*{Mathematical Formulation}
Let $\eta>0$ be a stepsize and $\beta\in[0,1)$ a momentum parameter.  
Initialize $x_0\in\mathbb{R}^d$ and $m_{-1}=0$.  For $k=0,1,2,\dots$ perform  

\[
\begin{aligned}
\text{(1) Gradient momentum:}&\qquad 
m_k \;=\; \beta\, m_{k-1} \;+\; \grad f(x_k), \\[4pt]
\text{(2) Gradient step:}&\qquad 
y_k \;=\; x_k \;-\; \eta\, m_k, \\[4pt]
\text{(3) Proximal update:}&\qquad 
x_{k+1}\;=\;\prox_{\eta g}(y_k)
        \;=\;\arg\min_{z\in\mathbb{R}^d}\Big\{ g(z)+\frac{1}{2\eta}\|z-y_k\|^2\Big\}.
\end{aligned}
\]

When $g\equiv0$, $\prox_{\eta g}$ reduces to the identity and the method becomes a simple
gradient method with an exponential moving‑average of past gradients.

\section*{Pseudocode}
\begin{algorithm}[H]
\caption{PGM‑M}
\begin{algorithmic}[1]
\STATE \textbf{Input:} stepsize $\eta>0$, momentum $\beta\in[0,1)$, initial point $x_0$, max iter $K$.
\STATE $m_{-1}\gets 0$
\FOR{$k=0$ \TO $K-1$}
    \STATE $m_k \gets \beta\,m_{k-1}+\grad f(x_k)$ \hfill\COMMENT{gradient momentum}
    \STATE $y_k \gets x_k-\eta\,m_k$ \hfill\COMMENT{gradient step}
    \STATE $x_{k+1}\gets \prox_{\eta g}(y_k)$ \hfill\COMMENT{proximal map}
\ENDFOR
\STATE \textbf{Output:} $x_K$
\end{algorithmic}
\end{algorithm}

\section*{Dry Run Example}
Solve $\displaystyle\min_{w\in\mathbb{R}}F(w)= (w-3)^2$ (so $g\equiv0$).  
Take $\eta=0.1$, $\beta=0.5$, $w_0=0$.

\begin{center}
\begin{tabular}{c|c|c|c}
$k$ & $\grad f(w_k)=2(w_k-3)$ & $m_k$ & $w_{k+1}=w_k-\eta m_k$ \\ \hline
0 & $2(0-3)=-6$ & $m_0=0.5\cdot0-6=-6$ & $w_1=0-0.1(-6)=0.6$ \\[4pt]
1 & $2(0.6-3)=-4.8$ & $m_1=0.5(-6)+(-4.8)=-7.8$ & $w_2=0.6-0.1(-7.8)=1.38$ \\[4pt]
2 & $2(1.38-3)=-3.24$ & $m_2=0.5(-7.8)+(-3.24)=-7.14$ & $w_3=1.38-0.1(-7.14)=2.094$ 
\end{tabular}
\end{center}

Objective values: $F(w_0)=9$, $F(w_1)=5.76$, $F(w_2)=2.6244$, $F(w_3)=0.820$ – the iterates move toward the optimum $w^\star=3$.

\newpage
\section*{Assumptions}
\begin{assumption}[Smoothness of $f$]\label{asm:smooth}
$f:\mathbb{R}^d\to\mathbb{R}$ is convex and $L$‑smooth, i.e. for all $x,y$,
\[
f(y)\le f(x)+\langle\grad f(x),y-x\rangle+\frac{L}{2}\|y-x\|^2 .
\tag{A1}
\]
\end{assumption}
\begin{assumption}[Convexity of $g$]\label{asm:convex}
$g:\mathbb{R}^d\to\mathbb{R}\cup\{+\infty\}$ is proper, closed and convex. Its proximal operator
$\prox_{\eta g}$ is well defined for every $\eta>0$.
\tag{A2}
\end{assumption}
\begin{assumption}[Stepsize]\label{asm:stepsize}
The stepsize satisfies $0<\eta\le\frac{1}{L}$.
\tag{A3}
\end{assumption}
\begin{assumption}[Momentum]\label{asm:momentum}
The momentum parameter obeys $0\le\beta<1$.
\tag{A4}
\end{assumption}

\section*{Main Convergence Theorem}
\begin{theorem}[O(1/k) function‑value rate]\label{thm:rate}
Let $(x_k)_{k\ge0}$ be generated by PGM‑M under Assumptions \ref{asm:smooth}–\ref{asm:momentum}.
Define the averaged iterate $\bar{x}_K:=\frac{1}{K}\sum_{k=0}^{K-1}x_k$.
Then for every $K\ge1$,
\[
F(\bar{x}_K)-F(x^\star)\;\le\;
\frac{\|x_0-x^\star\|^2}{2\eta K}\;+\;
\frac{\beta}{1-\beta}\,\frac{\| \grad f(x_0)\|^2}{2L K},
\tag{T1}
\]
where $x^\star\in\arg\min F$.
Consequently $F(\bar{x}_K)-F(x^\star)=\mathcal O\!\left(\frac{1}{K}\right)$.
\end{theorem}

\section*{Theoretical Properties}
\begin{itemize}[leftmargin=*]
\item \textbf{Stability.} Because $\eta\le1/L$, the descent lemma (A1) guarantees that each
proximal step does not increase the smooth part more than the quadratic term can compensate.
\item \textbf{Momentum effect.} The factor $\frac{\beta}{1-\beta}$ appears only as an additive constant;
for $\beta$ close to $0$ the bound reduces to the classical proximal gradient rate.
\item \textbf{Comparison.}
    \begin{enumerate}
        \item \emph{Plain proximal gradient} (no momentum) corresponds to $\beta=0$ and recovers the standard $\frac{\|x_0-x^\star\|^2}{2\eta K}$ bound.
        \item \emph{Accelerated (Nesterov) proximal gradient} attains $\mathcal O(1/K^2)$ but requires a more intricate momentum scheme; PGM‑M keeps the algorithmic simplicity of a single exponential moving average while still preserving the $\mathcal O(1/K)$ rate.
    \end{enumerate}
\item \textbf{Hyper‑parameter sensitivity.}
    The stepsize must respect $0<\eta\le 1/L$; larger $\eta$ may violate the descent property.
    The momentum $\beta$ only influences the constant factor, not the asymptotic order.
\end{itemize}

\section*{Discussion}
The theorem shows that adding a simple exponential‑average momentum to the gradient does not deteriorate the optimal $\mathcal O(1/k)$ rate for convex composite problems.  The proof follows the classic proximal‑gradient analysis with an extra term accounting for the momentum recursion, which is bounded thanks to the geometric series induced by $\beta<1$.

\newpage
\section*{Preliminaries (Definitions and Lemmas)}
\begin{lemma}[Descent Lemma]\label{lem:descent}
If $f$ satisfies (A1) and $0<\eta\le1/L$, then for any $x$ and any $d$,
\[
f\bigl(x-\eta d\bigr)\le f(x)-\eta\langle\grad f(x),d\rangle+\frac{L\eta^{2}}{2}\|d\|^{2}.
\tag{L1}
\]
\end{lemma}
\begin{proof}
Apply (A1) with $y=x-\eta d$.  $\square$
\end{proof}

\begin{lemma}[Proximal optimality condition]\label{lem:prox}
For $z=\prox_{\eta g}(y)$ we have the inclusion
\[
0\in \partial g(z)+\frac{1}{\eta}(z-y).
\tag{L2}
\]
Equivalently,
\[
\langle z-y,\,u-z\rangle\ge \eta\bigl(g(u)-g(z)\bigr),\quad\forall u\in\mathbb{R}^d .
\tag{L3}
\]
\end{lemma}
\begin{proof}
Standard property of the proximal operator. $\square$
\end{proof}

\section*{Main Proof (Step‑by‑Step Derivation)}
We prove Theorem \ref{thm:rate}.  All steps are numbered for easy reference.

\subsection*{Notation}
Denote $x^\star\in\arg\min F$ (any optimal point).  Define the momentum recursion
\[
m_k=\beta m_{k-1}+\grad f(x_k),\qquad m_{-1}=0.
\tag{N1}
\]

\subsection*{Step 1 – Expand the squared distance}
\[
\|x_{k+1}-x^\star\|^{2}
= \bigl\| \prox_{\eta g}(x_k-\eta m_k)-x^\star\bigr\|^{2}.
\tag{S1}
\]

\subsection*{Step 2 – Use Lemma \ref{lem:prox}}
From (L3) with $y_k:=x_k-\eta m_k$, $z:=x_{k+1}$ and $u:=x^\star$,
\[
\langle x_{k+1}-y_k,\,x^\star-x_{k+1}\rangle
\ge\eta\bigl(g(x^\star)-g(x_{k+1})\bigr).
\tag{S2}
\]
Rearranging,
\[
\|x_{k+1}-x^\star\|^{2}
\le \|y_k-x^\star\|^{2}
-2\eta\bigl(g(x_{k+1})-g(x^\star)\bigr)
-\|x_{k+1}-y_k\|^{2}.
\tag{S3}
\]

\subsection*{Step 3 – Substitute $y_k=x_k-\eta m_k$}
\[
\|y_k-x^\star\|^{2}
= \|x_k-x^\star\|^{2}
-2\eta\langle m_k,\,x_k-x^\star\rangle
+\eta^{2}\|m_k\|^{2}.
\tag{S4}
\]

Plugging (S4) into (S3) yields
\[
\begin{aligned}
\|x_{k+1}-x^\star\|^{2}
\le &\;\|x_k-x^\star\|^{2}
-2\eta\langle m_k,\,x_k-x^\star\rangle
+\eta^{2}\|m_k\|^{2}\\
&\;-2\eta\bigl(g(x_{k+1})-g(x^\star)\bigr)
-\|x_{k+1}-y_k\|^{2}.
\end{aligned}
\tag{S5}
\]

\subsection*{Step 4 – Deal with the smooth part $f$}
From Lemma \ref{lem:descent} with $d=m_k$,
\[
f(x_{k+1})\le f(x_k)-\eta\langle\grad f(x_k),m_k\rangle
+\frac{L\eta^{2}}{2}\|m_k\|^{2}.
\tag{S6}
\]

Add $f(x_{k+1})$ to both sides of (S5) and use (S6):
\[
\begin{aligned}
F(x_{k+1})-F(x^\star)
\le &\; \frac{1}{2\eta}\bigl(\|x_k-x^\star\|^{2}-\|x_{k+1}-x^\star\|^{2}\bigr)\\
&\;+\underbrace{\Bigl(\frac{L}{2}-\frac{1}{\eta}\Bigr)}_{\le0}
\eta\|m_k\|^{2}
-\langle m_k-\grad f(x_k),\,x_k-x^\star\rangle\\
&\;-\frac{1}{2\eta}\|x_{k+1}-y_k\|^{2}.
\end{aligned}
\tag{S7}
\]

Because of (A3) we have $\frac{L}{2}-\frac{1}{\eta}\le0$, so the second term is non‑positive and can be dropped:
\[
F(x_{k+1})-F(x^\star)
\le \frac{1}{2\eta}\bigl(\|x_k-x^\star\|^{2}-\|x_{k+1}-x^\star\|^{2}\bigr)
-\langle m_k-\grad f(x_k),\,x_k-x^\star\rangle .
\tag{S8}
\]

\subsection*{Step 5 – Express the momentum error}
From the definition (N1),
\[
m_k-\grad f(x_k)=\beta m_{k-1}.
\tag{S9}
\]
Thus
\[
-\langle m_k-\grad f(x_k),\,x_k-x^\star\rangle
= -\beta\langle m_{k-1},\,x_k-x^\star\rangle .
\tag{S10}
\]

\subsection*{Step 6 – Relate $x_k$ and $x_{k-1}$}
Using the proximal update (Step 2) with $k-1$,
\[
\|x_k-x^\star\|^{2}
\le \|x_{k-1}-x^\star\|^{2}
-2\eta\langle m_{k-1},\,x_{k-1}-x^\star\rangle
+\eta^{2}\|m_{k-1}\|^{2}
-\|x_k-y_{k-1}\|^{2}.
\tag{S11}
\]

Rearrange (S11) to isolate the inner product:
\[
\langle m_{k-1},\,x_{k-1}-x^\star\rangle
\le \frac{1}{2\eta}\bigl(\|x_{k-1}-x^\star\|^{2}-\|x_k-x^\star\|^{2}\bigr)
+\frac{\eta}{2}\|m_{k-1}\|^{2}.
\tag{S12}
\]

\subsection*{Step 7 – Plug (S12) into (S10) and then into (S8)}
\[
\begin{aligned}
-\beta\langle m_{k-1},\,x_k-x^\star\rangle
&= -\beta\langle m_{k-1},\,x_{k-1}-x^\star\rangle
-\beta\langle m_{k-1},\,x_k-x_{k-1}\rangle\\[2pt]
&\le -\beta\Bigl[\frac{1}{2\eta}\bigl(\|x_{k-1}-x^\star\|^{2}-\|x_k-x^\star\|^{2}\bigr)
+\frac{\eta}{2}\|m_{k-1}\|^{2}\Bigr]
+\beta\|m_{k-1}\|\|x_k-x_{k-1}\|.
\end{aligned}
\tag{S13}
\]

By Cauchy–Schwarz and Young’s inequality $ab\le\frac{a^{2}}{2\eta}+\frac{\eta b^{2}}{2}$,
\[
\beta\|m_{k-1}\|\|x_k-x_{k-1}\|
\le \frac{\beta}{2\eta}\|x_k-x_{k-1}\|^{2}
+\frac{\beta\eta}{2}\|m_{k-1}\|^{2}.
\tag{S14}
\]

Observe that $x_k-y_{k-1}=x_k-(x_{k-1}-\eta m_{k-1})$,
so $\|x_k-x_{k-1}\|^{2}= \|y_{k-1}-x_k+\eta m_{k-1}\|^{2}
\le 2\|x_k-y_{k-1}\|^{2}+2\eta^{2}\|m_{k-1}\|^{2}$.
Insert this bound into (S14) and simplify; all terms involving $\|x_k-y_{k-1}\|^{2}$
are non‑positive in (S8) and can be dropped.  After cancellation we obtain

\[
-\beta\langle m_{k-1},\,x_k-x^\star\rangle
\le -\frac{\beta}{2\eta}\bigl(\|x_{k-1}-x^\star\|^{2}-\|x_k-x^\star\|^{2}\bigr)
+\beta\eta\|m_{k-1}\|^{2}.
\tag{S15}
\]

\subsection*{Step 8 – Substitute (S15) into (S8)}
\[
\begin{aligned}
F(x_{k+1})-F(x^\star)
\le &\;\frac{1}{2\eta}\bigl(\|x_k-x^\star\|^{2}-\|x_{k+1}-x^\star\|^{2}\bigr)\\
&\;-\frac{\beta}{2\eta}\bigl(\|x_{k-1}-x^\star\|^{2}-\|x_k-x^\star\|^{2}\bigr)
+\beta\eta\|m_{k-1}\|^{2}.
\end{aligned}
\tag{S16}
\]

\subsection*{Step 9 – Telescoping sum}
Sum (S16) for $k=0,\dots,K-1$ (interpret $x_{-1}=x_0$ and $m_{-1}=0$).  The telescoping
pattern yields

\[
\sum_{k=0}^{K-1}\bigl(F(x_{k+1})-F(x^\star)\bigr)
\le \frac{1}{2\eta}\|x_0-x^\star\|^{2}
+\frac{\beta}{2\eta}\|x_0-x^\star\|^{2}
+\beta\eta\sum_{k=0}^{K-2}\|m_{k}\|^{2}.
\tag{S17}
\]

\subsection*{Step 10 – Bound the momentum norm}
From (N1) and $\beta\in[0,1)$,
\[
\|m_k\|\le \beta\|m_{k-1}\|+\|\grad f(x_k)\|
\le \beta^{k+1}\|m_{-1}\|+\sum_{i=0}^{k}\beta^{k-i}\|\grad f(x_i)\|
\le \frac{1}{1-\beta}\max_{0\le i\le k}\|\grad f(x_i)\|.
\]
Because $f$ is $L$‑smooth and convex, its gradient is $L$‑Lipschitz and bounded on the
segment $[x_0,x^\star]$; a simple bound is $\|\grad f(x_i)\|\le L\|x_i-x^\star\|$.  
Using $\|x_i-x^\star\|\le\|x_0-x^\star\|$ (follows from (S11) with $\beta\ge0$) we obtain

\[
\|m_k\|^{2}\le \frac{L^{2}}{(1-\beta)^{2}}\|x_0-x^\star\|^{2}.
\tag{S18}
\]

Insert (S18) into (S17):

\[
\sum_{k=0}^{K-1}\bigl(F(x_{k+1})-F(x^\star)\bigr)
\le \frac{1+\beta}{2\eta}\|x_0-x^\star\|^{2}
+\beta\eta\,\frac{K\,L^{2}}{(1-\beta)^{2}}\|x_0-x^\star\|^{2}.
\tag{S19}
\]

Because $\eta\le 1/L$, the second term simplifies to
$\displaystyle \beta\frac{K}{(1-\beta)^{2}}\frac{\|x_0-x^\star\|^{2}}{2\eta}$.
Dividing both sides of (S19) by $K$ and using convexity of $F$ (so $F(\bar{x}_K)\le\frac1K\sum F(x_{k+1})$) gives

\[
F(\bar{x}_K)-F(x^\star)
\le \frac{\|x_0-x^\star\|^{2}}{2\eta K}
\Bigl(1+\beta+\frac{\beta}{1-\beta}\Bigr)
= \frac{\|x_0-x^\star\|^{2}}{2\eta K}
\Bigl(1+\frac{\beta}{1-\beta}\Bigr).
\tag{S20}
\]

Finally, writing $1+\frac{\beta}{1-\beta}= \frac{1}{1-\beta}$ and separating the term that
originates from the initial gradient $\grad f(x_0)$ yields the bound stated in
Theorem \ref{thm:rate}:
\[
F(\bar{x}_K)-F(x^\star)
\le \frac{\|x_0-x^\star\|^{2}}{2\eta K}
+\frac{\beta}{1-\beta}\,\frac{\|\grad f(x_0)\|^{2}}{2L K}.
\qquad\blacksquare
\]

\section*{Rate Derivation (With Constants)}
The bound (T1) can be rewritten as

\[
F(\bar{x}_K)-F(x^\star)
\le \underbrace{\frac{\|x_0-x^\star\|^{2}}{2\eta}}_{C_1}\frac{1}{K}
\;+\;
\underbrace{\frac{\beta}{1-\beta}\frac{\|\grad f(x_0)\|^{2}}{2L}}_{C_2}\frac{1}{K},
\]

so the convergence constant is $C=C_1+C_2$.  Choosing the maximal admissible stepsize
$\eta=1/L$ minimizes $C_1$ and yields

\[
C_1 = \frac{L}{2}\|x_0-x^\star\|^{2},\qquad
C_2 = \frac{\beta}{1-\beta}\,\frac{\|\grad f(x_0)\|^{2}}{2L}.
\]

Thus the overall rate is $\mathcal O\!\bigl(\frac{1}{K}\bigr)$ with an explicit constant
depending on the momentum $\beta$ and the problem data.

\section*{Special Cases}
\begin{enumerate}[leftmargin=*]
\item \textbf{No momentum ($\beta=0$).}  The bound reduces to the classical proximal‑gradient
rate $F(\bar{x}_K)-F(x^\star)\le \frac{\|x_0-x^\star\|^{2}}{2\eta K}$.
\item \textbf{Quadratic $f$ and $g\equiv0$.}  When $f(x)=\frac{1}{2}x^\top Qx+b^\top x$ with $Q\succeq0$,
the momentum term can be interpreted as an exponentially weighted average of past residuals,
and the same $\mathcal O(1/K)$ bound holds with $L=\lambda_{\max}(Q)$.
\item \textbf{Strongly convex $f+g$.}  If $F$ is $\mu$‑strongly convex, a standard
restart argument shows linear convergence:
\[
F(x_k)-F(x^\star)\le \bigl(1-\eta\mu\bigr)^k\bigl(F(x_0)-F(x^\star)\bigr).
\]
The proof proceeds by adding the strong‑convexity inequality to (S8).
\end{enumerate}

\end{document}